\section{Threats to Validity}
\label{sec:threats}

Several factors limit the generalizability of these results. Internally, RQ4 and RQ5 rely on offline policy replay with a synthetic ambiguity injection manually relabeling the first suspicious window as unknown rather than naturally occurring ambiguous windows, and the bottleneck classification used to compute ground-truth labels is itself rule-based rather than independently verified. Externally, all experiments use a single model (\texttt{Qwen2.5-7B-Instruct}), a single serving framework (vLLM), a single profiler (Nsight Systems), and two GPU platforms (L4, A100); results may not transfer to other model architectures, serving stacks, or GPU generations. The finding that GPU-counter signals (\texttt{gpu\_util\_recovered}, \texttt{memory\_pressure\_low}) are inoperative is specific to saturated, high-concurrency serving and may not hold under lighter load, where utilization is more likely to dip below the recovery threshold. The RQ4 composite score uses a fixed weighting (40\% accuracy, 25\% premature-stop avoidance, 20\% re-escalation avoidance, 15\% cost-efficiency) that is a judgment call rather than an empirically derived tradeoff, and a different weighting could reorder policy rankings. Finally, repetition counts are modest (five microbenchmark repetitions, three vLLM seeds, eighteen replay streams), which limits the statistical power of the reported confidence intervals.

\section{Conclusion}

This paper presented an adaptive GPU tracing controller that resolves the tension between profiling cost and diagnostic confidence in LLM serving systems by collecting cheap evidence continuously and activating short, bounded Nsight Systems bursts only when needed, stopping as soon as the kernel-level diagnosis has stabilized. Across NVIDIA L4 and A100 GPUs, six vLLM serving scenarios, and three microbenchmark workloads, the controller reduced profiler duration by 18.9\%--48.5\% and captured kernel instances by 42.9\%--78.2\% relative to manual fixed-window profiling, while preserving a diagnosis match rate of 1.0 and introducing less than 1\% mean p95 latency overhead. Among the six stopping policies evaluated, Hybrid Stop and Counter-Recovery Stop eliminate re-escalation risk at modest additional cost, Stability Stop and Marginal Utility Stop offer a cheaper evidence-driven alternative, and a single Fixed Burst is unsafe under ambiguous first-window conditions; of the nine candidate stopping signals, only diagnosis stability reliably avoids firing on unconverged diagnoses. Together these results show that adaptive, evidence-driven control of heavy GPU profiling can substantially reduce tracing cost without sacrificing diagnostic accuracy, offering a practical alternative to today's all-or-nothing approach to GPU profiling in production LLM serving.
